\documentclass[dvipsnames]{article}

\usepackage{csquotes}

% Set page size and margins
% Replace `letterpaper' with`a4paper' for UK/EU standard size
\usepackage[letterpaper,top=2cm,bottom=2cm,left=3cm,right=3cm,marginparwidth=1.75cm]{geometry}

% Useful packages
\usepackage{amsmath}
\usepackage{graphicx}
\usepackage{minted2}
\usepackage{amsthm}
\usepackage{mathtools}
\usepackage[colorlinks=true, allcolors=blue]{hyperref}
\usepackage{fancyhdr}
\usepackage{enumerate}
\usepackage{tikz-cd}
\usepackage{mdframed}
\usepackage{multicol}

\setminted{frame=single}
\newmintinline[java]{java}{}

\usepackage{array}
\newcolumntype{L}[1]{>{\raggedright\let\newline\\\arraybackslash\hspace{0pt}}m{#1}}
\newcolumntype{C}[1]{>{\centering\let\newline\\\arraybackslash\hspace{0pt}}m{#1}}
\newcolumntype{R}[1]{>{\raggedleft\let\newline\\\arraybackslash\hspace{0pt}}m{#1}}

\DeclareMathOperator*{\argmin}{arg\,min}

\usepackage{tikz}
\usetikzlibrary{shapes}
\usetikzlibrary{arrows}

\tikzset{
  treenode/.style = {align=center, inner sep=0pt, text centered,
    font=\sffamily},
  avl_node/.style = {treenode, circle, black, draw=black,
    text width=1.5em, very thick},% arbre rouge noir, noeud rouge
  black_node/.style = {treenode, circle, white, font=\sffamily\bfseries, draw=black,
    fill=black, text width=1.5em},% arbre rouge noir, noeud noir
  red_node/.style = {treenode, circle, red, draw=red,
    text width=1.5em, very thick},% arbre rouge noir, noeud rouge
  nil_leaf/.style = {treenode, rectangle, draw=black, fill=black,
    minimum width=0.3em, minimum height=0.3em}% arbre rouge noir, nil
}

\newcommand{\haskell}[1]{\mintinline{haskell}{#1}}
\newcommand{\racket}[1]{\mintinline{racket}{#1}}
\newcommand{\prolog}[1]{\mintinline{prolog}{#1}}

\DeclareMathOperator{\xor}{\mathsf{xor}}
\DeclareMathOperator{\tru}{\mathsf{tru}}
\DeclareMathOperator{\fls}{\mathsf{fls}}
\DeclarePairedDelimiter{\ceil}{\lceil}{\rceil}

\begin{document}
\;

\section*{Topic 3. Problem set (theory)}
\setcounter{section}{3}

\thispagestyle{fancy}
\lhead{Solution authored by \textcolor{red}{Gabdulla Tukay \texttt{g.tukay@innopolis.university}}}
\cfoot{Data Structures and Algorithms (Spring 2026), Innopolis University}

\theoremstyle{definition}
\newtheorem{problem}{Problem}[section]

\begin{problem}
  Determine the worst-case time complexity of the function \texttt{process} (in terms of the size of the input linked list --- $n$):

  \begin{minipage}{0.5\textwidth}
  \setlength{\leftmargini}{0em}
  \begin{enumerate}
    \item Write down the recurrence relation for $T(n)$ (the running time of the function \texttt{process}).
    \item Find the asymptotic complexity $T(n)$ using
    \begin{itemize}
      \item the substitution method \cite[\S 4.3]{CLRS-2022},
      \item the recursion tree method \cite[\S 4.4]{CLRS-2022}, or
      \item the master method \cite[\S 4.5]{CLRS-2022}.
    \end{itemize}
    Your solution should explicitly work through each step of the method, and for asymptotic proofs you must use formal definitions
    or known properties with explicit references (to a specific location in the book \cite{CLRS-2022} or a specific slide in the lecture).
  \end{enumerate}
\end{minipage}
\begin{minipage}{0.05\textwidth}
  \;
\end{minipage}
\begin{minipage}{0.45\textwidth}

  \begin{minted}[texcomments,escapeinside=||,frame=single,fontsize=\footnotesize,linenos]{javascript}
/* L is a linked list of size n */
function process(L)
  if L.is_empty()
    return 0
  else if L.size() == 1
    return L.head.value
  else
    lists = new array of 3 empty linked lists
    i = 0
    A = L.head
    while A is not None
      if (i mod 4) > 0
        lists[i mod 4].add(A.value)
      A = A.next
      i = i + 1
    result = process(lists[1])
    result += process(lists[2])
    result += process(lists[3])
    return result
  \end{minted}
\end{minipage}

%%%%%%%%%%%%%%%%%%%%%%%%%%%%%%%%%%%%%%%%%%%%%%%%%
% Start of answer and full solution
%%%%%%%%%%%%%%%%%%%%%%%%%%%%%%%%%%%%%%%%%%%%%%%%%
\color{purple}

\textbf{(a) Recurrence relation:}

\begin{mdframed}
  \color{purple}
  \emph{Replace with your recurrence relation $T(n) = \ldots$}
\end{mdframed}

\textbf{(b) Asymptotic complexity:}

\begin{mdframed}
  \color{purple}
  \textbf{Method used:} \emph{Substitution method / Recursion tree method / Master method}
  
  \textbf{Asymptotic complexity:} $T(n) = \Theta(\ldots)$
  
  \begin{proof}[\textbf{Solution steps}]
    \emph{Replace with your step-by-step solution using the chosen method.}
    \emph{Include explicit references to CLRS or lecture slides where needed.}
  \end{proof}
\end{mdframed}
%%%%%%%%%%%%%%%%%%%%%%%%%%%%%%%%%%%%%%%%%%%%%%%%%
% End of answer and full solution
%%%%%%%%%%%%%%%%%%%%%%%%%%%%%%%%%%%%%%%%%%%%%%%%%

\end{problem}

\clearpage
\begin{problem}
  Solve the following recurrence relations using the master method (see \cite[Theorem~4.1]{CLRS-2022}), if applicable.
  For each recurrence relation, your solution should contain:
  \begin{itemize}
    \item The answer --- the asymptotic complexity for $T(n)$ expressed in $\Theta$-notation, or an indication that the master method is not applicable.
    \item If the master method is applicable, then:
    \begin{itemize}
      \item Indicate which case of the master theorem is used (Case 1, 2, or 3).
      \item Explicitly write down the conditions that need to be checked for the chosen case (specialized for this particular recurrence).
      \item Prove that the conditions are satisfied. For asymptotic proofs, you must use formal definitions or known properties with explicit references (to a specific location in the book \cite{CLRS-2022} or a specific slide in the lecture).
    \end{itemize}
    \item If the master method is not applicable, then explicitly justify why.
  \end{itemize}
  \begin{enumerate}
    \item $T(n) = 4 T(n / 16) + \sqrt{n} \cdot \log_2 n$
    \item $T(n) = 5 T(4 n / 9) + n^2$
    \item $T(n) = 9 T(n / 4) + \log_2 \left(n!\right)$
    \item $T(n) = 3 T(\sqrt[3]{n}) + n$               
    \item $T(n) = \frac{1}{2026} T(2026 n) + n \log_{2026} n$
  \end{enumerate}

%%%%%%%%%%%%%%%%%%%%%%%%%%%%%%%%%%%%%%%%%%%%%%%%%
% Start of answer and full solution
%%%%%%%%%%%%%%%%%%%%%%%%%%%%%%%%%%%%%%%%%%%%%%%%%
\color{purple}

\begin{enumerate}
  \item \textcolor{black}{$T(n) = 4 T(n / 16) + \sqrt{n} \cdot \log_2 n$}
  \begin{mdframed}
    \color{purple}
    \textbf{Answer:} $T(n) = \Theta(\ldots)$ \emph{or Master method is not applicable.}
    
    \begin{proof}[\textbf{Solution}]
      \emph{Replace with your solution. If applicable, indicate the case, write down conditions, and prove them.}
      \emph{If not applicable, justify why.}
    \end{proof}
  \end{mdframed}
  
  \item \textcolor{black}{$T(n) = 5 T(4 n / 9) + n^2$}
  \begin{mdframed}
    \color{purple}
    \textbf{Answer:} $T(n) = \Theta(\ldots)$ \emph{or Master method is not applicable.}
    
    \begin{proof}[\textbf{Solution}]
      \emph{Replace with your solution. If applicable, indicate the case, write down conditions, and prove them.}
      \emph{If not applicable, justify why.}
    \end{proof}
  \end{mdframed}
  
  \item \textcolor{black}{$T(n) = 9 T(n / 4) + \log_2 \left(n!\right)$}
  \begin{mdframed}
    \color{purple}
    \textbf{Answer:} $T(n) = \Theta(\ldots)$ \emph{or Master method is not applicable.}
    
    \begin{proof}[\textbf{Solution}]
      \emph{Replace with your solution. If applicable, indicate the case, write down conditions, and prove them.}
      \emph{If not applicable, justify why.}
    \end{proof}
  \end{mdframed}
  
  \item \textcolor{black}{$T(n) = 3 T(\sqrt[3]{n}) + n$}
  \begin{mdframed}
    \color{purple}
    \textbf{Answer:} $T(n) = \Theta(\ldots)$ \emph{or Master method is not applicable.}
    
    \begin{proof}[\textbf{Solution}]
      \emph{Replace with your solution. If applicable, indicate the case, write down conditions, and prove them.}
      \emph{If not applicable, justify why.}
    \end{proof}
  \end{mdframed}
  
  \item \textcolor{black}{$T(n) = \frac{1}{2026} T(2026 n) + n \log_{2026} n$}
  \begin{mdframed}
    \color{purple}
    \textbf{Answer:} $T(n) = \Theta(\ldots)$ \emph{or Master method is not applicable.}
    
    \begin{proof}[\textbf{Solution}]
      \emph{Replace with your solution. If applicable, indicate the case, write down conditions, and prove them.}
      \emph{If not applicable, justify why.}
    \end{proof}
  \end{mdframed}
\end{enumerate}
%%%%%%%%%%%%%%%%%%%%%%%%%%%%%%%%%%%%%%%%%%%%%%%%%
% End of answer and full solution
%%%%%%%%%%%%%%%%%%%%%%%%%%%%%%%%%%%%%%%%%%%%%%%%%

\end{problem}

\clearpage
\begin{problem}
  Consider a modification of the merge sort algorithm \cite[\S 2.3]{CLRS-2022}
  that stops recursive division on arrays of size $k$ or less ($k \leq n$).
  On an array of size $k$, the ``conquer'' phase consists of sorting this array
  using insertion sort \cite[\S 2.2]{CLRS-2022}.
  
  \begin{enumerate}
    \item Write down the pseudocode of the modified sorting algorithm described above.
    \item What is the time complexity of the modified sorting algorithm in the worst case (in terms of $n$ and $k$)?
      Give your answer using $\Theta$-notation and provide justification for your estimate.
    \item What is the time complexity of the modified sorting algorithm in the best case (in terms of $n$ and $k$)?
      Give your answer using $\Theta$-notation and provide justification for your estimate.
  \end{enumerate}

%%%%%%%%%%%%%%%%%%%%%%%%%%%%%%%%%%%%%%%%%%%%%%%%%
% Start of answer and full solution
%%%%%%%%%%%%%%%%%%%%%%%%%%%%%%%%%%%%%%%%%%%%%%%%%
\color{purple}

\textbf{(a) Pseudocode:}

\begin{proof}[\textbf{Pseudocode}]
  \emph{Replace with your pseudocode.}
\end{proof}

\textbf{(b) Worst case time complexity:}

\begin{mdframed}
  \color{purple}
  \textbf{Time complexity:} $T(n) = \Theta(\ldots)$
  
  \begin{proof}[\textbf{Justification}]
    \emph{Replace with your justification.}
  \end{proof}
\end{mdframed}

\textbf{(c) Best case time complexity:}

\begin{mdframed}
  \color{purple}
  \textbf{Time complexity:} $T(n) = \Theta(\ldots)$
  
  \begin{proof}[\textbf{Justification}]
    \emph{Replace with your justification.}
  \end{proof}
\end{mdframed}
%%%%%%%%%%%%%%%%%%%%%%%%%%%%%%%%%%%%%%%%%%%%%%%%%
% End of answer and full solution
%%%%%%%%%%%%%%%%%%%%%%%%%%%%%%%%%%%%%%%%%%%%%%%%%

\end{problem}

\begin{thebibliography}{CormenLeiserson09}

\bibitem[CLRS]{CLRS-2022}
Cormen, T.H., Leiserson, C.E., Rivest, R.L. and Stein, C., 2022. \emph{Introduction to algorithms, Fourth Edition}. MIT press.

\end{thebibliography}

\end{document}
